\documentclass[10pt]{article}
\usepackage{icmlko}

\icmlmeta{IP 파운데이션 월드모델}{502}{자물쇠는 사슬이었다}{표적 표를 5분의 1로 잠근 것은 한 줄이었다. 그것을 풀자 스텝 501 의 「못 가른다」가 자료가 아니라 부트 설계의 성질로 드러났다}

\begin{document}

\twocolumn[
\icmlheader

\begin{abstract}
\textbf{이 논문은 두 가지를 한다.} 첫째, 세 사이클 동안 「표적 표가 464행에서 안 자란다」던
진단의 원인을 찾았다 --- \textbf{원료는 처음부터 표 옆에 있었다.} 위키 문서가 해결된 개체는
\textbf{3{,}933} 인데 수확기의 표적 선택이 「판 유보」로 묶여 있어 \textbf{815(20.7\,\%)}만
시야에 들었고, 수확기는 그 815 중 810 을 \emph{이미} 다 돌았다. 잠금을 풀어 3{,}896 개체를
받으니(성공률 99.06\,\%) \textbf{팔 B 가 464 $\to$ 2{,}050 으로 4.42배} 자랐다. 층 ③(무관
배경)의 이득은 그 4.4배 표에서 \textbf{자리를 안 옮겼고}($+$0.059774 $\to$ $+$0.058238)
SE 는 $1/\sqrt{n}$ 이 예측하는 대로 \textbf{2.23배 줄었다}(0.0226 $\to$ 0.010137).
둘째, \textbf{스텝 501 의 네 번째 주장을 철회한다.} 501 은 도메인 안 $\Delta\rho=+$0.087763
에 SE 0.044957 을 붙여 「0 과 못 가른다」고 적었다. 그 부트스트랩은 재표집마다 $N_{\min}$
걸러내기를 \emph{다시} 적용해 \textbf{추정량 자체를 흔든다}. 원표본의 묶음 집합에 조건부로
바꾸면 \textbf{씨앗 24개 전부에서 $T$ 를 넘고}, 501 의 설계로는 \textbf{24개 중 13개}만
넘는다 --- \textbf{동전 던지기였다}. 다만 그 안쪽 몫 43.62\,\% 자체가 \textbf{칸막이의
성질}이라 개체 규모 분위로 자르면 8.96\,\%, 액션 시각 분위로 자르면 21.04\,\% 다.
\end{abstract}
]

\section{잠금은 한 줄이었다}

수확기 \verb|targets()| 의 첫 줄이 \verb|hold = json.loads(HOLD.read_text())["유보키"]|
였다. 그래서 표적이 \textbf{판 유보}로 묶였다.

\begin{center}\small
\begin{tabular}{lr}
\hline
단계 & 수 \\
\hline
캐시 파일 & 10{,}960 \\
\textbf{위키 문서가 해결된 개체} & \textbf{3{,}933} \\
\textbf{수확기가 볼 수 있던 것} & \textbf{815 (20.7\,\%)} \\
그 815 중 이미 받아 둔 것 & 810 \\
\hline
\end{tabular}
\end{center}

\textbf{수집이 게을렀던 것이 아니다 --- 시야가 5분의 1이었다.} 그리고 같은 파일의
독스트링이 「옛 자(판 유보 3{,}775 에 몇 행 붙나)는 이 사이클의 자가 아니다」라고 적어
놓고 \textbf{표적 선택에는 그 옛 자를 남겼다}. 선언과 배선이 어긋난 자리다.

\claimx{표적 표가 안 자란 원인은 자료의 부족이 아니라 \textbf{표적 선택의 범위}다.
잠금을 풀면 위키 삼중쌍이 584 $\to$ 2{,}570, 팔 B 가 464 $\to$ 2{,}050(4.42배)이 된다.}
{새 표에서 층 ③의 $\Delta\rho$ 부호가 바뀌면 「자란 표가 더 나은 표」라는 전제가 죽는다.
실측은 $+$0.058238 로 부호가 그대로다.}

\section{층 ③은 4.4배 표에서 자리를 안 옮긴다}

\begin{center}\small
\begin{tabular}{lrr}
\hline
 & 팔 B $=$ 464 & 팔 B $=$ 2{,}050 \\
\hline
$\Delta\rho$(① $\to$ ①$+$③) & $+$0.059774 & $+$0.058238 \\
SE(군집 부트) & 0.0226 & \textbf{0.010137} \\
문턱 $T$ & 0.0452 & \textbf{0.020275} \\
판정 & 든다 & \textbf{든다} \\
\hline
\end{tabular}
\end{center}

점추정이 2.6\,\% 움직이는 동안 SE 가 2.23배 줄었다. $1/\sqrt{4.42}=0.476$ 이 예측하는
값과 사실상 같다. \textbf{다만 이 비교는 깨끗하지 않다} --- 표가 커지면 층 ③의 배경 풀도
같이 커지므로 「표본이 커진 것」과 「배경이 좋아진 것」을 이 실험은 못 가른다. 그 한계를
측정 전에 사전등록에 적었다.

\section{스텝 501 의 「못 가른다」를 철회한다}

501 은 \textbf{$\Delta\rho_{\text{within}}=+$0.087763} 에 SE $0.044957$ · $T=0.089915$ 를
붙여 \textbf{「못 가른다」}고 적었다. \textbf{그 SE 는 정정 대상인 죽은 숫자다} ---
값을 부풀린 것이 아니라 \textbf{추정량이 흔들린 것}이다. 501 의 부트는 재표집한 표본 위에서
$N_{\min}$ 걸러내기를 다시 불렀고, 도서 20 · 모바일 24 · 웹툰 24 · 팝업 24 가 문턱 20
\emph{바로 위}에 앉아 있어 묶음 집합이 재표집마다 4$\sim$8 로 출렁였다. SE 는 「그 추정량의
표집 변동」이어야 하는데 그것은 \textbf{두 잡음의 합}이다.

\begin{center}\small
\begin{tabular}{lrr}
\hline
부트 설계 & $T$ 를 넘는 씨앗 & 평균 SE \\
\hline
재적용 (501 이 쓴 것) & \textbf{13 / 24} & 0.043829 \\
원표본 조건부 (정본) & \textbf{24 / 24} & \textbf{0.040672} \\
\hline
\end{tabular}
\end{center}

씨앗 평균 95\,\% 구간은 $[+0.006467,\;+0.166147]$ 이고 \textbf{0을 품는 씨앗이 0개}다.
$N_{\min}\in\{1,5,10,15,20,25\}$ 를 두 설계에 대고 얻은 \textbf{12칸 중 붉은 칸은 하나}
뿐이다 --- \textbf{(재적용, $N_{\min}=20$)}, 곧 501 이 고른 단 하나의 조합이다.
같은 러너가 501 의 코드를 부르지 않고 자를 다시 구현해 501 의 수를 먼저 재현했다:
$\Delta\rho_{\text{within}}$ 차 $4.69\times10^{-7}$, 안쪽 몫 차 $4.53\times10^{-8}$.

\claimx{501 의 네 번째 주장(「그 $+$0.087763 도 0 과 못 가른다」)을 \textbf{철회한다}.
부트 설계를 추정량이 흔들리지 않게 바꾸면 씨앗 24개 전부에서 $T$ 를 넘는다. 501 의
「못 가른다」는 자료의 성질이 아니라 \textbf{설계의 성질}이었다.}
{원표본 조건부 설계에서 씨앗 비율이 95\,\% 밑으로 내려가면 죽는다. 그리고 이산 판정
지점이 부트 잡음과 같은 크기이면 판정을 내지 않고 구간만 싣기로 미리 등록했다 --- 이번엔
비율이 100\,\% 라 그 칸을 안 썼다.}

\section{그러나 「안쪽 몫」은 칸막이의 성질이다}

501 이 실은 「안쪽 몫 43.62\,\%」는 \textbf{도메인으로 잘랐을 때}의 수다. 같은 팔·같은
$\Delta\rho$ 에 칸막이만 바꾸면:

\begin{center}\small
\begin{tabular}{lrr}
\hline
칸막이 & 안쪽 짝 & 안쪽 몫 \\
\hline
도메인 (501 의 것) & 28{,}951 & 43.62\,\% \\
개체 규모 분위 4 & 26{,}667 & \textbf{8.96\,\%} \\
액션 시각 분위 4 & 26{,}681 & \textbf{21.04\,\%} \\
\hline
\end{tabular}
\end{center}

최대$-$최소가 \textbf{34.66\,\%p} 다. 칸막이 정의는 값을 보기 전에 사전등록에 못 박았다.

\claimx{「이득이 안쪽인가 바깥쪽인가」는 자료의 성질이 아니라 \textbf{칸막이의 성질}이다.
따라서 501 의 「안쪽 몫 43.62\,\%」는 \textbf{「도메인으로 잘랐을 때」라는 조건 없이는 못
선다}. within/between 을 쓸 때는 어느 칸막이인지를 수와 같은 줄에 적어야 한다.}
{세 칸막이의 안쪽 몫이 15\,\%p 안에 모이면 죽는다. 실측 34.66\,\%p.}

\section{자물쇠는 사슬이었다}

물리적 장소가 있는 도메인(팝업·시장팝업)의 행은 \textbf{59 $\to$ 118 로 정확히 두 배}가
됐다. 그런데 물리 장을 붙이는 팔 A 는 \textbf{39 $\to$ 38} 로 자라지 않았다 ---
새로 들어온 118행 중 \textbf{75행에 좌표가 없다}. 지오코딩이 옛 유보 집합에만 돌았기
때문이다. 그 75를 붙이면 팔 A 상한이 118 이고 $T$ 가 $0.2983 \to 0.169$ 로 내려가
\textbf{처음으로 판정 가능 구간}에 든다.

\claimx{병목은 하나가 아니라 사슬이다. 표적 선택 자물쇠를 풀자 \textbf{다음 자물쇠가
좌표라는 것이 수로 드러났다}(75행). 이것은 새 원천 0으로 풀 수 있다.}
{좌표 75행을 붙였는데도 팔 A 가 118 에 못 미치거나 $T$ 가 0.25 아래로 안 내려가면 죽는다.}

\section{틀린 예측을 신고한다}

사전등록한 아홉 예측 중 \textbf{여섯이 맞고 셋 중 둘이 이 사이클의 표적이었다}. 표적이던
Q5(「팔 A $\ge$ 78」)와 Q6(「팔 A 의 $T<0.25$」)이 \textbf{둘 다 틀렸다}. 그리고 \textbf{틀린
이유를 같은 문서가 이미 적고 있었다} --- 사전등록 ⓪절이 「팔 A 의 천장은 44이고 그 42가
이미 서울이다」라고 세어 놓았는데, Q5 의 문턱은 「팝업 개체가 두 배가 되니 팔 A 도 두 배」
라는 다른 전제로 세웠다. \textbf{한 문서 안의 두 문장이 어긋나 있었고 아무도 대조하지
않았다.} 이 병은 스텝 501 이 낳은 노트가 자기 반증조건에서 저지른 것과 같은 병이다.

\end{document}
